\section{Konsep dan Definisi Kunci}
\label{sec:konsep-kunci}

Bagian ini menetapkan istilah yang dipakai sepanjang buku.

\subsection{Instruction Set Architecture}

\begin{definisi}[Instruction Set Architecture]
\term{Instruction Set Architecture} (ISA) adalah kumpulan instruksi
yang dapat dijalankan sebuah prosesor beserta aturan penggunaannya:
register, mode pengalamatan, dan mekanisme interupsi.
\end{definisi}

ISA adalah kontrak antara perangkat lunak dan perangkat keras. Selama
program hanya memakai instruksi yang ada pada ISA, program itu dapat
dijalankan pada prosesor yang mengimplementasikan ISA tersebut.

\subsection{Instruksi, Operand, dan Mnemonic}

\begin{definisi}[Instruksi dan operand]
\term{Instruksi} adalah perintah yang dijalankan prosesor.
\term{Operand} adalah data atau tempat data yang dikenai instruksi,
baik berupa nilai, register, maupun alamat memori.
\end{definisi}

\begin{definisi}[Mnemonic]
\term{Mnemonic} adalah nama pendek yang mewakili sebuah instruksi
dalam bahasa assembly, misalnya \asm{MOV} untuk memindahkan data dan
\asm{ADD} untuk menjumlahkan.
\end{definisi}

Contohnya, pada \lstasm|ADD AX, 5|, \asm{ADD} adalah mnemonic
instruksi, sedangkan \reg{AX} dan 5 adalah operand.

\subsection{Register}

\begin{definisi}[Register]
\term{Register} adalah tempat penyimpanan kecil (8 atau 16 bit) yang
berada di dalam prosesor. Akses ke register jauh lebih cepat daripada
akses ke memori.
\end{definisi}

Jumlah register 8086 terbatas, sehingga programmer harus
mengelolanya secara sadar. Pembahasan lengkap ada pada bab tentang
register dan \eng{flag}.

\subsection{Siklus Instruksi}

\begin{definisi}[Siklus instruksi]
\term{Siklus instruksi} adalah rangkaian \eng{fetch}, \eng{decode},
dan \eng{execute} untuk satu instruksi.
\end{definisi}

Satu instruksi dapat melibatkan beberapa kali pengaksesan bus.
Instruksi yang mengakses memori umumnya lebih mahal daripada
instruksi yang hanya bekerja pada register.

\subsection{Memory-Mapped I/O dan Port I/O}

\begin{definisi}[Memory-mapped I/O dan port I/O]
\term{Memory-mapped I/O} adalah teknik menghubungkan perangkat
masukan/keluaran ke ruang alamat memori. \term{Port I/O} adalah
teknik yang memakai ruang alamat terpisah dan instruksi khusus
masukan/keluaran.
\end{definisi}

Kedua teknik dipakai pada keluarga 8086. Area video, misalnya, dapat
diakses melalui alamat memori, sedangkan perangkat lain diakses
melalui port. Pembahasan operasionalnya ada pada bab tentang
interupsi BIOS/DOS.

\subsection{Ringkasan Definisi}

\begin{table}[H]
    \centering
    \caption{Ringkasan istilah kunci dan tempat pembahasannya.}
    \label{tab:istilah-kunci}
    \small
    \begin{tabular}{@{}p{3.6cm}p{5.8cm}p{4.0cm}@{}}
        \toprule
        \textbf{Istilah} & \textbf{Pengertian singkat} &
        \textbf{Dibahas di} \\
        \midrule
        ISA & Kontrak instruksi antara perangkat lunak dan keras &
        Bab tentang arsitektur 8086 \\
        \addlinespace
        Operand & Data atau tempat data yang dikenai instruksi &
        Bab tentang mode pengalamatan \\
        \addlinespace
        Register & Penyimpan kecil di dalam prosesor &
        Bab tentang register \\
        \addlinespace
        Siklus instruksi & \eng{Fetch}, \eng{decode},
        \eng{execute} & Bab tentang arsitektur 8086 \\
        \addlinespace
        Segmentasi & Pembagian memori menjadi segmen 64 KB &
        Bab tentang segmentasi \\
        \addlinespace
        Interupsi & Mekanisme memanggil layanan sistem &
        Bab tentang interupsi BIOS/DOS \\
        \addlinespace
        \eng{Stack} & Penyimpan sementara dengan aturan LIFO &
        Bab tentang stack dan prosedur \\
        \addlinespace
        Array dan string & Kumpulan data berurutan di memori &
        Bab tentang array dan string \\
        \bottomrule
    \end{tabular}
\end{table}
